\subsection*{2. Compression vs. Low-Pass Filtering}

Compression and low-pass filtering both work in the frequency domain but pick their coefficients differently. Compression just ranks everything by magnitude and keeps whatever's strongest, doesn't matter what frequency it's at. Low-pass filtering instead keeps everything inside some fixed radius and zeros out the rest, even if that discarded stuff actually had a lot of energy in it.

Natural images tend to keep their important edge info in the high frequencies. Since compression keeps whatever has the highest magnitude, it ends up preserving those edges anyway. At threshold $0.001$ we only kept $0.1\%$ of the frequencies and the dog was still recognizable, just blocky. Low-pass filtering doesn't care about that though, it just chops high frequencies no matter what. Even at radius $10$, which keeps more coefficients, the cat image ended up so blurred that all the fur texture was basically gone.

So overall: compression adapts to the actual content, so the image stays structurally sharp but gets blocky. Low-pass filtering just grabs a fixed band of frequencies, so the image gets more and more blurry as the radius shrinks.
